\chapter{Manual de usuario}

\section*{Nota previa}

Este manual documenta el uso del programa \texttt{octrees-benchmark-coordenadas-polares}, que parte
de la implementación base del Octree lineal desarrollada por Viñambres et al. \cite{Vinambres}. El presente trabajo se centra en el
diseño e implementación de las técnicas de poda a nivel de nodo hoja
sobre las estructuras de datos existentes (\textit{Octrees}) y sus correspondientes algoritmos de búsqueda.

\section{Instalación}

\subsection{Dependencias}

\paragraph{LASTools}

En primer lugar es necesario instalar las dependencias de LASTools, listadas
en \url{https://github.com/LAStools/LAStools}:

\begin{verbatim}
sudo apt-get install libjpeg62 libpng-dev libtiff-dev libjpeg-dev \
    libz-dev libproj-dev liblzma-dev libjbig-dev libzstd-dev \
    libgeotiff-dev libwebp-dev liblzma-dev libsqlite3-dev
\end{verbatim}

A continuación, se clona el repositorio y se compila:

\begin{verbatim}
git clone --depth 1 https://github.com/LAStools/LAStools lib/LAStools
(cd lib/LAStools && cmake -B build -DCMAKE_BUILD_TYPE=Release \
    -DCMAKE_INSTALL_PREFIX=$PWD/../LASlib . && \
    cmake --build build -- -j && cmake --install build)
rm -rf lib/LAStools
\end{verbatim}

\paragraph{PCL versión 1.15 (Opcional)}

Se debe obtener el código fuente de la versión 1.15 desde
\url{https://github.com/PointCloudLibrary/pcl/releases} y compilarlo. La
carpeta de instalación usada por defecto es \texttt{\textasciitilde/local/pcl},
aunque puede modificarse, junto con la versión buscada, editando el
fichero \texttt{CMakeLibraries.cmake}.

\begin{verbatim}
wget https://github.com/PointCloudLibrary/pcl/releases/download/\
pcl-1.15.0/source.tar.gz
tar xvf source.tar.gz
cd pcl && mkdir build && cd build
cmake .. -DCMAKE_INSTALL_PREFIX=$HOME/local/pcl -DCMAKE_BUILD_TYPE=Release
make -j2
make -j2 install
\end{verbatim}

Si PCL no se encuentra durante la compilación, el proyecto se compilará
igualmente con normalidad, pero sin soporte para los algoritmos de
\textit{benchmark} basados en el Octree y el KD-Tree de PCL.

\paragraph{PAPI (Opcional, para perfilado de caché)}

\begin{verbatim}
wget https://github.com/icl-utk-edu/papi/releases/download/\
papi-7-2-0-t/papi-7.2.0.tar.gz -P lib
tar xvf lib/papi-7.2.0.tar.gz -C lib && rm lib/papi-7.2.0.tar.gz \
    && mv lib/papi-7.2.0 lib/papi/
(cd lib/papi/src && ./configure --prefix=$(pwd)/.. && make -j \
    && make install)
\end{verbatim}

\subsection{Compilación}

Una vez instaladas las dependencias necesarias, dentro del directorio del
proyecto basta con ejecutar:

\begin{verbatim}
cmake -B build -DCMAKE_BUILD_TYPE=Release .
cmake --build build
\end{verbatim}

Esto genera el ejecutable en \texttt{build/octrees-benchmark}.

Para más comodidad, existe el script de bash \texttt{compile.sh}, que ejecuta esos comandos además de limpiar el directorio de build (recomendado borrar el direcorio siempre que se compile después de introducir cambios en el código).

\subsection{Verificación de la instalación: Tests}

El proyecto incluye una batería de tests unitarios mediante GoogleTest, que
permite comprobar que la instalación y compilación se realizaron
correctamente. Para habilitarla, se debe compilar el proyecto con la opción
\texttt{-DBUILD\_TESTS=ON}:

\begin{verbatim}
cmake -B build -DCMAKE_BUILD_TYPE=Release -DBUILD_TESTS=ON .
cmake --build build
\end{verbatim}

Tras compilar el proyecto, todos los tests se pueden ejecutar mediante:

\begin{verbatim}
make                # compila la biblioteca y todos los ejecutables de test
ctest --output-on-failure
\end{verbatim}

También es posible ejecutar cada batería de tests de forma individual, por
ejemplo:

\begin{verbatim}
./tests/test_points
./tests/test_encoders
./tests/test_octree
./tests/test_octree_advanced
\end{verbatim}

Para ejecutar un caso de test concreto con CTest, se utiliza la opción
\texttt{-R} seguida del nombre del test, por ejemplo:

\begin{verbatim}
ctest -R LinearOctreeTest.RadiusSearch
\end{verbatim}

Adicionalmente, bajo \texttt{tests/example.cpp} se incluye un pequeño ejemplo
de uso de las funcionalidades de la biblioteca, que se compila
automáticamente junto con el proyecto principal:

\begin{verbatim}
make test_library
./tests/test_library
\end{verbatim}

\section{Ejecución}

El programa se ejecuta a través del binario generado en la compilación,
\texttt{build/octrees-benchmark}, configurando su comportamiento mediante
parámetros de línea de comandos. A continuación se muestra un ejemplo de
invocación, correspondiente a uno de los experimentos descritos en esta
memoria: una búsqueda completa secuencial de vecinos
sobre la nube \texttt{Paris\_Luxembourg\_6}, comparando los distintos reordenamientos existentes sobre los dos algorimtmos de búsqueda sobre los que se trabajó para una selección de radios.

\begin{verbatim}
./build/octrees-benchmark --kernels "sphere" -i "data/paris_lille/Lille_0.las" \
    -o "out_tfg_reorders_v3/full" -e "hilb" --kernels "cube,sphere" -r "2.0,5.0,8.0" -s "all" --sequential --repeats 1 -a "neighborsPrune,neighborsStruct" --local-reorder "none,polar,cartesian" --max-leaf "256" --umbral-poda "64"
\end{verbatim}

Para reproducir el conjunto completo de experimentos descritos en los
capítulos de Pruebas y Discusión de Resultados de esta memoria, el
proyecto incluye una serie de scripts:

\begin{itemize}
    \item \texttt{bench\_neighbors\_reorders.bash}: búsquedas de vecinos de centros aleatorios.
    \item \texttt{bench\_neighbors\_full\_reorders.bash}: búsquedas de vecinos de coberturas completas.
    \item \texttt{bench\_neighbors\_parallel.bash}: búsquedas de vecinos para diferentes escalas de paralelización, tanto de centros aleatorios como coberturas completas.
    \item \texttt{bench\_neighbors\_debug\_range\_selector.bash}: búsquedas reducidas (60 búsquedas), sobre el mismo conjunto de nubes, con la opción de depuración de podas habilitadas, devolviendo archivos con el número de puntos podados por hoja. (Modificar con precaución: generan archivos muy pesados, escribiendo una línea por hoja evaluada en cada búsqueda).
    \item Adicionalmente, se encuentran los scripts del estudio de partida: \item \texttt{bench\_locality.bash}, \item \texttt{bench\_memory.bash}, \item \texttt{bench\_neighbors.bash}, \item \texttt{bench\_neighbors\_temp.bash}
\end{itemize}

Asimismo, bajo la carpeta \texttt{plots\_reorders} se incluyen todos los scripts y notebooks en Python empleados para generar las figuras presentadas en esta memoria. Se puede consultar cada gráfica para cada una de las nubes de puntos estudiadas, ya que no todas fueron incluidas en la memoria. En la carpeta \texttt{plots} se encuentran las figuras de los resultados publicados por Viñambres et al., también disponibles para consulta.



\subsection{Formato de entrada}

El programa toma como entrada nubes de puntos en formato binario
\texttt{.LAS}, gestionadas internamente mediante la biblioteca LASTools.

\subsection{Formato de salida}

Los resultados de cada ejecución se almacenan en la ruta indicada mediante el parámetro \texttt{-o}/\texttt{--output}. El programa genera ficheros con las métricas temporales y de efectividad de poda correspondientes a la
configuración de parámetros especificada, que sirven de entrada a los scripts de generación de gráficas.

\subsection{Ejecución de la verificación de funcionamiento de las podas}
En la ruta \texttt{test\_funcionalidad\_podas/test\_vecindades.ipynb} está el test de Python que comprueba si los resultados usando reordenamientos son idénticos a los resultados sin hacer uso de la optimización (función \texttt{verificar\_integridad\_vecinos}). Los parámetros a introducir son los siguientes:
\begin{itemize}
    \item \textbf{data\_path}: Ruta de los archivos \texttt{.csv} devueltos por las ejecuciones del programa.
    \item \textbf{clouds}: El test comprobará todos los casos posibles solamente para las nubes aquí especificadas.
\end{itemize}

\section{Parámetros del programa} 

A continuación se detallan todos los parámetros aceptados por el programa. Para consultar esta misma información desde la propia terminal, se puede invocar la ayuda integrada mediante:

\begin{verbatim}
./build/octrees-benchmark --help
\end{verbatim}

\renewcommand{\arraystretch}{1.3}
\begin{longtable}{p{1.6cm} p{3.3cm} p{8.8cm}}
\caption{Parámetros principales del programa \texttt{octrees-benchmark}.} \\
\toprule
\textbf{Opción} & \textbf{Alias} & \textbf{Descripción} \\
\midrule
\endfirsthead
\toprule
\textbf{Opción} & \textbf{Alias} & \textbf{Descripción} \\
\midrule
\endhead
\bottomrule
\endfoot

\texttt{-h} & \texttt{--help} & Muestra el mensaje de ayuda. \\

\texttt{-i} & \texttt{--input} & Ruta del fichero de entrada. \\

\texttt{-c} & \texttt{--container-type} &
Tipo de contenedor a utilizar. Por defecto: \texttt{AoS}. Valores posibles:
\texttt{SoA}, \texttt{AoS}. \\

\texttt{-o} & \texttt{--output} & Ruta del fichero de salida. \\

\texttt{-r} & \texttt{--radii} &
Radios de búsqueda a evaluar (separados por comas, p. ej.\
\texttt{2.5,5.0,7.5}). \\

\texttt{-v} & \texttt{--kvalues} &
Valores de $k$ para las búsquedas KNN (separados por comas, p. ej.\
\texttt{10,50,250,1000}). \\

\texttt{-s} & \texttt{--searches} &
Número de búsquedas a realizar (con centros aleatorios, salvo que se
indique \texttt{--sequential}). Si se indica \texttt{all}, se realiza una
búsqueda sobre la totalidad de la nube (con indexación secuencial). \\

\texttt{-t} & \texttt{--repeats} &
Número de repeticiones a realizar para cada configuración del
\textit{benchmark}. \\

\texttt{-k} & \texttt{--kernels} &
Especifica los kernels de búsqueda a utilizar (separados por comas o
\texttt{all}). Valores posibles: \texttt{sphere}, \texttt{cube},
\texttt{square}, \texttt{circle}. \\

\texttt{-a} & \texttt{--search-algo} &
Especifica los algoritmos de búsqueda a ejecutar (separados por comas o
\texttt{all}). Por defecto:
\texttt{neighborsPtr,neighbors,neighborsPrune,neighborsStruct}.
Valores posibles: \newline
\textbf{Búsqueda por radio:} \newline
\hspace*{0.4cm}\textbullet\ \texttt{neighborsPtr} -- búsqueda básica sobre
el Octree de punteros \newline
\hspace*{0.4cm}\textbullet\ \texttt{neighbors} -- búsqueda básica sobre el
Octree lineal \newline
\hspace*{0.4cm}\textbullet\ \texttt{neighborsPrune} -- búsqueda optimizada
sobre el Octree lineal con poda de octantes \newline
\hspace*{0.4cm}\textbullet\ \texttt{neighborsStruct} -- búsqueda optimizada
mediante rangos de índices \newline
\hspace*{0.4cm}\textbullet\ \texttt{neighborsApprox} -- búsqueda aproximada
(cotas superior/inferior), requiere \texttt{--approx-tol} \newline
\hspace*{0.4cm}\textbullet\ \texttt{neighborsUnibn} -- búsqueda mediante
unibnOctree \newline
\hspace*{0.4cm}\textbullet\ \texttt{neighborsPCLKD} -- búsqueda mediante el
KD-Tree de PCL (si está disponible) \newline
\hspace*{0.4cm}\textbullet\ \texttt{neighborsPCLOct} -- búsqueda mediante el
Octree de PCL (si está disponible) \newline
\hspace*{0.4cm}\textbullet\ \texttt{neighborsPico} -- búsqueda mediante
PicoTree \newline
\textbf{Búsqueda KNN:} \newline
\hspace*{0.4cm}\textbullet\ \texttt{KNNV2} -- búsquedas KNN sobre el Octree
lineal \newline
\hspace*{0.4cm}\textbullet\ \texttt{KNNNanoflann} -- búsquedas KNN mediante
nanoflann \newline
\hspace*{0.4cm}\textbullet\ \texttt{KNNPCLKD} -- búsqueda KNN mediante el
KD-Tree de PCL (si está disponible) \newline
\hspace*{0.4cm}\textbullet\ \texttt{KNNPCLOCT} -- búsqueda KNN mediante el
Octree de PCL (si está disponible) \newline
\hspace*{0.4cm}\textbullet\ \texttt{KNNPico} -- búsqueda KNN mediante
PicoTree \\

\texttt{-e} & \texttt{--encodings} &
Selecciona las codificaciones SFC con las que reordenar la nube antes de
las búsquedas (separadas por comas o \texttt{all}). Por defecto:
\texttt{all}. Valores posibles: \newline
\hspace*{0.4cm}\textbullet\ \texttt{none} -- sin codificación; el Octree
lineal no se construirá con esta opción \newline
\hspace*{0.4cm}\textbullet\ \texttt{mort} -- reordenamiento mediante SFC de
Morton \newline
\hspace*{0.4cm}\textbullet\ \texttt{hilb} -- reordenamiento mediante SFC de
Hilbert \\

\texttt{--} & \texttt{--local-reorders} &
Selecciona los reordenamientos locales a nivel de hoja a aplicar sobre la nube antes de las búsquedas (separados por comas o \texttt{all}). Por defecto:
\texttt{all}. Valores posibles: \newline
\hspace*{0.4cm}\textbullet\ \texttt{none} -- sin reordenamiento local \newline
\hspace*{0.4cm}\textbullet\ \texttt{polar} -- reordenamiento local polar
(según el ángulo azimutal $\varphi$) \newline
\hspace*{0.4cm}\textbullet\ \texttt{cartesian} -- reordenamiento local
cartesiano (según las coordenadas $x$, $y$, $z$) \\

\texttt{-u} & \texttt{--umbral-poda} &
Umbral mínimo de puntos
(\texttt{threshold}) necesario para emplear el selector de rango y podar la
hoja. Por defecto: 16. Toda hoja con un número de puntos inferior a este
umbral se recorrerá de forma secuencial, sin aplicar la poda. \\

\texttt{--} & \texttt{--debug-ranges} &
Activa el registro
(\textit{logging}) detallado de la selección de rangos en cada hoja
conflictiva visitada durante las búsquedas. \\

\texttt{--} & \texttt{--debug} &
Activa el modo de depuración general (mide los tiempos de construcción del
Octree y de codificación). \\

\texttt{--} & \texttt{--build-enc} &
Ejecuta \textit{benchmarks} para la codificación y construcción de las
estructuras seleccionadas (aquellas con representante en \texttt{-a}\,/\,
\texttt{--search-algo}). \\

\texttt{--} & \texttt{--memory} &
Ejecuta un \textit{benchmark} simple para medir la memoria consumida por
una estructura, mediante perfilado del \textit{heap}. Valores posibles:
\texttt{ptrOct}, \texttt{linOct}, \texttt{unibnOct}, \texttt{nanoKD},
\texttt{pclOct}, \texttt{pclKD}, \texttt{picoTree}. \\

\texttt{--} & \texttt{--locality} &
Ejecuta \textit{benchmarks} para analizar la localidad de la nube de puntos
tras aplicar los reordenamientos indicados. \\

\texttt{--} & \texttt{--cache-profiling} &
Activa el perfilado de la memoria caché durante la ejecución de los
algoritmos de búsqueda, mediante PAPI. \\

\texttt{--} & \texttt{--check} &
Activa la comprobación de resultados (opción heredada; se recomienda usar
\texttt{avg\_result\_size} para verificar la corrección de los
resultados). \\

\texttt{--} & \texttt{--no-warmup} &
Desactiva la fase de calentamiento (\textit{warmup}) previa a las
mediciones. \\

\texttt{--} & \texttt{--approx-tol} &
Valores de tolerancia para la búsqueda aproximada (separados por comas,
p. ej.\ \texttt{10.0,50.0,100.0}). \\

\texttt{--} & \texttt{--num-threads} &
Lista de números de hilos a evaluar en el test de escalabilidad (separados
por comas, p. ej.\ \texttt{1,2,4,8,16,32}). Si no se especifica, OpenMP
emplea por defecto el número máximo de hilos disponibles y no se ejecuta
ningún test de escalabilidad. \\

\texttt{--} & \texttt{--sequential} &
Hace que el conjunto de búsquedas sea secuencial en lugar de aleatorio
(habitualmente más rápido). Se activa automáticamente al usar
\texttt{-s all}. \\

\texttt{--} & \texttt{--max-leaf} &
Número máximo de puntos por hoja del Octree (por defecto = 128). No aplica
al Octree de PCL. \\

\texttt{--} & \texttt{--pcl-oct- resolution} &
Tamaño mínimo de octante para la subdivisión en el Octree de PCL. \\

\end{longtable}

\section{Mensajes de error y problemas comunes}

A continuación se recogen algunas situaciones de error habituales y cómo
resolverlas:

\begin{itemize}
    \item \textbf{El programa no encuentra el fichero de entrada}: comprobar
    que la ruta indicada en \texttt{-i}/\texttt{--input} es correcta y que el
    fichero existe, tiene formato \texttt{.LAS} y tiene permisos de lectura.

    \item \textbf{Uso de 
\end{itemize}